\documentclass[12pt]{article}
\usepackage[utf8]{inputenc}
\usepackage[T1]{fontenc}
\usepackage{hyperref}
\title{Wali Language Resources for K-12 Curriculum Development}
\author{Compiled by Claude Sonnet 5.5 (claude-sonnet-5-5) for Andr\'as Kornai}
\date{}
\begin{document}
\maketitle

\section{Introduction}
Wali (ISO 639-3: wlx; Glottocode wali1263) is a Gur language spoken mainly in and around Wa in Ghana's Upper West Region, and also in the Northern Region. Wikipedia gives about 85,000 speakers (2013). It is written in Latin script, and three competing orthographies exist: from the Ghana Institute of Languages, GILLBT, and Baptist Mid-Missions. See \url{https://en.wikipedia.org/wiki/Wali_language_(Gur)}.

\section{Primary Educational Resources}
The searches found no Wali primers or textbooks.

\section{Linguistic and Reference Materials}
No Wali-specific dictionary or grammar turned up. The search returned \emph{A Dictionary and Grammatical Sketch of Dagaare} (Ali, Grimm and Bodomo, Language Science Press, 2021) at \url{https://library.oapen.org/handle/20.500.12657/48439}. Dagaare is a different language, so it is not a Wali resource and is not listed in the JSON. The only Wali source found is the Wikipedia overview.

\section{Religious and Community Resources}
None found in the searches.

\section{Digital and Audio Resources}
None found in the searches.

\section{Curriculum Development Context}
The orthography is not standardised, which is the first problem for any curriculum. A new search aimed at GILLBT, the Ghana Institute of Languages and SIL materials would be the next step. Little else is documented.
\end{document}